\subsection*{Objects, models, structures}

A. Mathematical objects and structures. From antiquity to the nineteenth century there was common agreement on the principal objects of study of a mathematician, namely those mentioned by Plato in the passage quoted earlier: numbers, magnitudes, and figures. Although, for the Greeks, this list should be enlarged to include the objects and phenomena proper to mechanics, astronomy, optics, and music, they nevertheless always made a clear distinction between these ``mathematical'' disciplines and arithmetic and geometry; and from the time of the Renaissance the former soon acquired the status of independent sciences.

Whatever the philosophical nuances which coloured the conception of mathematical objects in the mind of this or that mathematician or philosopher, there was at least one point of unanimous agreement: that these objects are given and that it does not lie in our power to ascribe arbitrary properties to them, any more than a physicist can alter a natural phenomenon. In truth, these views are partly the product of psychological reactions, which it is not for us to go into, but which every mathematician is well aware of when he exhausts himself in vain efforts to obtain a proof which seems perpetually to escape him. From there it is only one step to assimilating this resistance to external obstacles; and even today more than one mathematician, who parades an intransigent formalism, would voluntarily subscribe in his conscience to this avowal by Hermite: ``I believe that the numbers and functions of Analysis are not the arbitrary product of our minds; that they exist outside of us, with the same character of necessity as the objects of the material world, and we meet or discover them, and study them, just like the physicists, chemists, and zoologists ({[}27{]}, vol.~II, p.~398).

In the classical conception, mathematics was the study of only numbers and figures. But this official doctrine, to which every mathematician felt bound to give his verbal adherence, gradually became more and more of an intolerable constraint under the pressure of new ideas. The embarrassment of algebraists in the presence of negative numbers vanished only when analytical geometry provided them with a convenient ``interpretation''. But even in the eighteenth century d'Alembert (although a convinced ``positivist''), when discussing the question in the Encyclopédie {[}13{]}, suddenly lost courage after a column of somewhat confused explanations, and contented himself with concluding that ``the rules of algebraic operations on negative quantities are generally admitted by everyone and are generally received as correct, whatever interpretation is to be attached to these quantities''. As regards imaginary numbers, the scandal was far worse still; for if they were ``impossible'' roots and if (until about 1800) there was no known way of ``interpreting'' them, how could these undefinable objects be talked about without contradiction, and above all how could they be introduced? On this topic d'Alembert prudently kept his silence, and did not even state the questions, no doubt because he realized that he could not answer in terms other than those naively used by A. Girard a century earlier {[}9{]} : ``On pourroit dire : à quoy sert ces solutions qui sont impossibles ? Je réponds : pour trois choses, pour la certitude de la reigle générale, et qu'il n'y a point d'autres solutions, et pour son utilité''.

In analysis, the situation in the eighteenth century was hardly better. It was a happy coincidence that analytical geometry appeared, as if on cue, to provide an ``interpretation'', in the form of a geometrical figure, of the great creation of the 17th century, the notion of function, and thus to assist powerfully (in the hands of Fermat, Pascal, and Barrow) at the birth of infinitesimal calculus. But the philosophico-mathematical controversies to which the notions of infinitely small and indivisible quantities gave rise are only too well known. And although d'Alembert was on firmer ground here, and realized that in the ``metaphysics'' of the infinitesimal calculus there is nothing but the concept of limit, he was no more able than his contemporaries to comprehend the true meaning of expansions in divergent series, or to explain the paradox of correct results obtained by calculations with expressions which are devoid of all numerical interpretation. Finally, even in the domain of ``geometrical certainty'', the Euclidean façade was cracking. When Stirling, in 1717, did not hesitate to say that a certain curve has an ``imaginary double point at infinity'' (*), he would have been hard put to relate such an ``object'' to commonly received notions; and Poncelet, who, at the beginning of the nineteenth century, gave a considerable impetus to such ideas by founding projective geometry, was content to invoke as justification an entirely metaphysical ``principle of continuity''.

In these circumstances (and, paradoxically, at the very time when the ``absolute truth'' of mathematics was being proclaimed with the greatest insistence), the notion of proof seemed to become more and more blurred as the eighteenth century ran its course, because the notions employed, and their fundamental properties, could not be definitively fixed in the manner of the Greeks. The return to rigour, which was set in motion at the beginning of the nineteenth century, brought some improvements to this state of affairs, but did not correspondingly stem the flood of new ideas. Thus in algebra appeared the imaginaries of Galois, the ideal numbers of Kummer, followed by vectors and quaternions, n-dimensional spaces, multivectors, and tensors, not to mention Boolean algebra. Undoubtedly one of the great advances (which made possible a return to rigour without abandonment of any part of the conquests of the past) was the possibility of constructing ``models'' of these new notions in more classical terms. Thus ideal numbers and the imaginary numbers of Galois were interpreted in terms of the theory of congruences, n-dimensional geometry could be regarded (if one so wished) purely as a language for expressing the results of algebra ``of n variables''; and, as regards the classical imaginary numbers --- whose geometrical representation by the points of a plane marked the beginning of this efflorescence of algebra --- there was soon the choice between this geometrical ``model'' and an interpretation in terms of congruences. But mathematicians at last began to feel that this makeshift approach was contrary to the natural development of their science and that it should be legitimate, in mathematics, to work with objects which have no concrete ``interpretation''. ``It is not of the essence of mathematics'', said Boole in 1854, ``to be conversant with the ideas of number and quantity'' ({[}16 b{]}, p.~13) (**). The same considerations led Grassmann, in his ``Ausdehnungslehre'' of 1844, to present his calculus in a form in which the notions of number and geometrical object were completely excluded from the start (*). And a little later, Riemann, in his inaugural lecture, was careful not to speak of ``points'' but of ``determinations'' (Bestimmungsweise) in his description of ``manifolds n times extended'', and emphasized that in such a manifold the ``metrical relations'' (Massverhältnisse) ``can be studied only for abstract quantities and can be represented only by formulae; under certain assumptions they can however be decomposed into relations, each of which separately is capable of a geometrical representation, and thereby the results of calculations may be expressed in geometrical terms'' ({[}19{]}, p.~276).

From this moment onward, the enlargement of the axiomatic method was an accepted fact. Although for some time yet it was felt to be useful to check ``abstract'' results, whenever possible, against geometrical intuition, at least it was agreed that the ``classical'' objects were not the only legitimate objects of study for a mathematician. Precisely because of the multiplicity of possible ``interpretations'' or ``models'', it came to be recognized that the ``nature'' of mathematical objects is ultimately of secondary importance, and that it matters little, for example, whether a result is presented as a theorem of ``pure'' geometry or as a theorem of algebra via analytical geometry. In other words, the essence of mathematics --- that fleeting and insubstantial notion which hitherto could only be expressed by vague names such as ``reigle generale'' or ``metaphysic''---appeared as the study of relations between objects which do not of themselves intrude on our consciousness, but are known to us by means of some of their properties, namely those which serve as the axioms at the basis of their theory. Boole had understood this clearly in 1847 when he wrote that mathematics is concerned with ``operations considered in themselves, independently of the various ways in which they may be applied'' ({[}16 a{]}, p.~3). Hankel, in 1867, inaugurating the axiomatization of algebra, defended mathematics as being ``purely intellectual, a pure theory of forms whose purpose is the study, not of the combination of magnitudes or of their images, namely numbers, but objects of thought (Gedankendinge), which may correspond to concrete objects or relations, although such a correspondence is not necessary'' ({[}20{]}, p.~10). Cantor, in 1883, echoes this claim for a ``free mathematics'', proclaiming that ``mathematics is entirely free in its development, and its concepts are bound only by the necessity that they should be non-contradictory and coordinated with concepts previously introduced by precise definitions'' ({[}25{]}, p.~182). Finally, the revision of Euclidean geometry helped to disseminate and popularize these ideas. Pasch himself, although still attached to a certain ``reality'' of geometrical objects, realized that geometry is a fact independent of their signification, and consists purely in the study of their relations ({[}28{]}, p.~90): a conception which Hilbert pushed to its logical conclusion by emphasizing that even the names of the basic notions of a mathematical theory may be chosen at random (*), and which Poincaré expressed by saying that axioms are ``disguised definitions'', thereby completely reversing the point of view of the scholastic philosophers.

It is therefore tempting to assert that the modern notion of ``structure'' was substantially in existence by 1900, but in fact another thirty years of preparation were required before it made its full-fledged appearance. Certainly it is not difficult to recognize structures of the same species when they are of a sufficiently simple nature; with group-structures, for example, this point was attained in the middle of the nineteenth century. But at the same period Hankel was still struggling, without complete success, to extract the general ideas of a field and an extension field, which he managed to express only in the semi-metaphysical form of a ``principle of permanence'' {[}20{]}, and which were definitively formulated only by Steinitz forty years later. It has been especially difficult to escape from the feeling that mathematical objects are ``given'' together with their structure, and it has taken many years of functional analysis to make modern mathematicians familiar with the idea that, for example, there are several ``natural'' topologies on the set of rational numbers, and several measures on the real line. With this dissociation the passage to the general definition of structures such as has been presented in this book, was finally achieved.

B. Models and isomorphisms. In several places we have had occasion to refer to the notion of a ``model'' or an ``interpretation'' of a mathematical theory constructed with the help of another theory. This is not a recent idea, and it should undoubtedly be seen as a continually recurring manifestation of a deep-lying feeling of the unity of the various ``mathematical sciences''. If the traditional maxim ``All things are numbers'' of the early Pythagoreans is taken as authentic, it may be considered as the vestige of a first attempt to reduce the geometry and algebra of the time to arithmetic. Although the discovery of irrationals appeared to close this path for good, the reaction it provoked in Greek mathematics was a second attempt at synthesis, this time taking geometry as the basis and embracing, among other things, the methods of solution of algebraic equations inherited from the Babylonians (*). This conception reigned supreme until the fundamental reform of R. Bombelli and Descartes, which assimilated every measurement of magnitudes to a measurement of length (in other words, to a real number). But with the creation of analytical geometry by Descartes and Fermat, the trend was again reversed, and a far closer fusion of geometry and algebra was obtained, this time to the benefit of algebra. Descartes went further and at one stroke conceived of the essential unity of ``all the sciences which are commonly called mathematical\ldots{} Although their objects are different'', he wrote, ``they are all concordant with one another in that they are concerned only with the various ratios or proportions which are present'' ({[}10{]}, vol.~VI, pp.~19-20) (†). However, this point of view tended only to make algebra the fundamental mathematical science: a conclusion against which vigorous protests were uttered by Leibniz, who, as we have seen, had also conceived of a ``universal mathematics'', but on a far vaster scale and already in close accord with modern ideas. Making quite explicit the ``concordance'' of which Descartes spoke, he perceived for the first time what is effectively the general notion of isomorphism (which he called ``similitude'') and the possibility of ``identifying'' isomorphic relations or operations; as an example he cites addition and multiplication ({[}12 bis{]}, pp.~301-303). But these bold conceptions found no echo among his contemporaries, and the realization of his dreams had to await the expansion of algebra in the middle of the nineteenth century. We have already observed that it was at this point that the ``models'' multiplied and the habit of passing from one theory to another by simple changes of language was formed; the most striking example of the latter is perhaps that of duality in projective geometry, where the practice of printing theorems ``dual'' to each other side by side in two columns undoubtedly had a large share in the acceptance of the notion of isomorphism. On a more technical level, it is certain that the notion of isomorphic groups was known to Gauss for commutative groups, and to Galois for groups of permutations; it was in defined for general arbitrary groups about the middle of the nineteenth century (*). Thereafter, with each new axiomatic theory it became natural to define a notion of isomorphism; but it is only with the modern notion of structure that it has been finally recognized that every structure contains in itself a notion of isomorphism, and that it is superfluous to give a special definition for each species of structures.

C. The arithmetization of classical mathematics. The increasingly widespread use of the notion of ``model'' was also to permit the nineteenth century mathematicians to achieve the unification of mathematics dreamed of by the Pythagoreans. At the beginning of the century, whole numbers and continuous quantities appeared to be as irreconcilable as they were in antiquity; the real numbers were related to the notion of geometric magnitude (especially that of length), and the ``models'' of negative numbers and imaginary numbers were constructed on this basis. Even the rational numbers were traditionally associated with the idea of ``subdivision'' of a magnitude into equal parts. Only the integers remained apart, as ``exclusive products of our intellect'', as Gauss said in 1832, when contrasting them with the notion of space ({[}14{]}, vol.~VIII, p.~201). The first attempts, by Martin Ohm (1822), to bring together arithmetic and analysis were concerned with the rational numbers (positive and negative); they were taken up again around 1860 by several authors, notably Grassmann,

Hankel, and Weierstrass (in unpublished lectures). To Weierstrass is apparently due the idea of obtaining a ``model'' of the positive or negative rational numbers by considering classes of ordered pairs of natural integers. But the most important step still remained to be taken, namely to construct a ``model'' of the irrational numbers from within the theory of rational numbers. By 1870 this had become an urgent problem because of the necessity, after the discovery of ``pathological'' phenomena in analysis, to purge every trace of geometrical intuition and the vague notion of ``magnitude'' from the definition of real numbers. The problem was in fact solved, at about this time, almost simultaneously by Cantor, Dedekind, Méray, and Weierstrass, using quite different methods from one another.

From then on the integers became the foundation of all classical mathematics. Furthermore, the ``models'' founded on arithmetic acquired still greater importance with the extension of the axiomatic method and the conception of mathematical objects as free creations of the human intellect. For there remained a limitation to the freedom claimed by Cantor, namely the limitation raised by the question of ``existence'', which had already occupied the Greeks, and which arose now so much more urgently precisely because all appeal to intuitive representation was now abandoned. We shall see later what a philosophico-mathematical maelström was to be generated by the notion of ``existence'' in the early years of the twentieth century. But in the nineteenth century this stage had not yet been reached, and to prove the existence of a mathematical object having a given set of properties meant simply, just as for Euclid, to ``construct'' an object with the required properties. This was precisely the purpose of the arithmetical ``models''; once the real numbers had been ``interpreted'' in terms of integers, then the same was done for complex numbers and Euclidean geometry, thanks to analytic geometry, and likewise for all the new algebraic objects introduced since the beginning of the century. Finally, in a discovery which achieved great fame, Beltrami and Klein obtained Euclidean ``models'' of the non-Euclidean geometries of Lobatschevsky and Riemann and thereby ``arithmetized'' (and completely justified) these theories which at first had aroused so much distrust.

D. The axiomatization of arithmetic. It was natural, in this line of development, that attention should next be directed toward the foundations of arithmetic itself, and in fact this is what happened around 1880. Apparently, before the nineteenth century, no one had attempted to define addition and multiplication in any other way than by a direct appeal to intuition. Leibniz alone, faithful to his principles, pointed out explicitly that ``truths'' as ``obvious'' as \(2 + 2 = 4\) are no less capable of proof, if one reflects on the definitions of the numbers which appear therein ({[}12 b{]}, vol.~IV, p.~403; cf.~{[}12 bis{]}, p.~203); and he did not by any means regard the commutativity of addition and multiplication as self-evident (*). He did not, however, push his reflections on this subject any further, and in the middle of the nineteenth century progress had still not been made in this direction. Even Weierstrass, whose lectures contributed greatly to the spread of ``arithmetizing'' point of view, seems not to have realized the necessity of a logical clarification of the theory of integers. The first steps in this direction are apparently due to Grassmann who, in 1861 ({[}17{]}, vol.~II₂, p.~295), gave a definition of addition and multiplication of integers and proved their fundamental properties (commutativity, associativity, distributivity), using nothing but the operation \(x \rightarrow x + 1\) and the principle of induction. The latter had been clearly conceived and used for the first time in the seventeenth century by B. Pascal ({[}11{]}, vol.~III, p.~456) (†) --- though more or less explicit applications of it are to be found in the mathematics of antiquity --- and it was in current use by mathematicians from the second half of the seventeenth century. But it was not until 1888 that Dedekind ({[}26{]}, vol.~III, pp.~359-361) enunciated a complete system of axioms for arithmetic (reproduced three years later by Peano and usually known under his name {[}29 a{]}), which contained in particular a precise formulation of the principle of induction (which Grassmann had used without enunciating it explicitly).

With this axiomatization it seemed that the definitive foundations of mathematics had been attained. But, in fact, at the very moment when the axioms of arithmetic were being clearly formulated, arithmetic itself was being dethroned from this role of primordial science in the eyes of many mathematicians (beginning with Dedekind and Peano), in favour of the most recent of mathematical theories, namely the theory of sets; and the controversies which were to surround the notion of integers cannot be isolated from the great ``crisis of foundations'' of the years 1900-1930.

